\chapter{Leitura de Sensores Analógicos (ADC)}
\label{cap:adc}

Muitos sensores — potenciômetros, sensores de luminosidade (LDR), sensores de temperatura analógicos — não respondem apenas com ``ligado'' ou ``desligado'': eles produzem uma \textbf{tensão contínua e variável}, proporcional à grandeza física medida. Para ler esse tipo de sinal, usamos um \textbf{ADC} (\emph{Analog-to-Digital Converter}, ou conversor analógico-digital) — a porta de entrada de praticamente todo projeto de monitoramento IoT.

\section{Como um ADC funciona}

Um ADC amostra uma tensão analógica de entrada e a converte em um número inteiro, dentro de uma faixa determinada pela sua \textbf{resolução}, medida em bits. O \ESP{} possui ADCs de \textbf{12 bits}, o que significa que cada leitura retorna um valor entre 0 e $2^{12}-1 = 4095$ — quanto mais próxima a tensão de entrada estiver da tensão máxima suportada, mais próximo de 4095 será o valor lido.

\begin{atencao}
Se você já programou com placas Arduino clássicas (como o Arduino Uno), tenha atenção: lá, \code{analogRead()} retorna valores de \textbf{10 bits} (0 a 1023). No \ESP{}, por padrão, a resolução é de \textbf{12 bits} (0 a 4095) — o dobro de precisão, mas também uma faixa de valores diferente da que você talvez esteja acostumado.
\end{atencao}

\begin{esp32box}
O \ESP{} possui duas unidades ADC independentes: \textbf{ADC1} (8 canais, nos pinos GPIO32 a GPIO39) e \textbf{ADC2} (10 canais, compartilhados com o rádio Wi-Fi — por isso, leituras em ADC2 podem falhar ou ficar imprecisas enquanto o Wi-Fi está ativo). Por essa razão, esta apostila usa exclusivamente pinos da \textbf{ADC1} nos exemplos a seguir.
\end{esp32box}

\begin{table}[h]
\centering
\begin{tabular}{@{}cc@{}}
\toprule
\textbf{Pino GPIO} & \textbf{Unidade ADC} \\
\midrule
GPIO32 & ADC1 \\
GPIO33 & ADC1 \\
GPIO34 & ADC1 \\
GPIO35 & ADC1 \\
GPIO36 & ADC1 \\
GPIO39 & ADC1 \\
\bottomrule
\end{tabular}
\caption{Alguns pinos da ADC1 disponíveis no ESP32.}
\end{table}

\section{Lendo um sensor analógico}

No framework Arduino, uma única função basta para realizar uma leitura: \code{analogRead(pino)}.

\begin{codigoc}{src/main.cpp — lendo um potenciômetro ou LDR}
#define PINO_SENSOR 34

void setup() {
    Serial.begin(115200);
}

void loop() {
    int leitura_bruta = analogRead(PINO_SENSOR);
    float tensao = (leitura_bruta / 4095.0f) * 3.3f;

    Serial.print("Bruto: ");
    Serial.print(leitura_bruta);
    Serial.print("  |  Aproximadamente: ");
    Serial.print(tensao);
    Serial.println(" V");

    delay(500);
}
\end{codigoc}

\begin{atencao}
A conversão \code{(leitura\_bruta / 4095.0f) * 3.3f} usada acima é uma \textbf{aproximação} — assume uma relação perfeitamente linear entre o valor lido e a tensão real, o que não é exatamente verdade na prática devido a não linearidades do hardware do ADC. Para a maioria dos sensores simples em um projeto de aprendizado — um LDR, um potenciômetro, um sensor de temperatura analógico — essa aproximação linear costuma ser suficiente.
\end{atencao}

\section{De leitura bruta a informação: um sensor de luminosidade}

Uma leitura analógica isolada raramente é o objetivo final — o valor bruto do ADC quase sempre precisa ser \textbf{traduzido} para uma grandeza com significado físico, e depois \textbf{registrado} ou \textbf{avaliado}. O exemplo a seguir lê um LDR (fotorresistor), estima um percentual de luminosidade e usa o próprio monitor serial como ferramenta de depuração e acompanhamento — a técnica de \emph{print debugging} apresentada no \cref{cap:entrada-saida}, agora aplicada a um sensor real:

\begin{codigoc}{Estimando luminosidade e registrando via Serial}
#define PINO_LDR 34

void setup() {
    Serial.begin(115200);
}

void loop() {
    int leitura_bruta = analogRead(PINO_LDR);
    float luminosidade_pct = (leitura_bruta / 4095.0f) * 100.0f;

    Serial.print("[SENSOR] LDR bruto=");
    Serial.print(leitura_bruta);
    Serial.print(" luminosidade=");
    Serial.print(luminosidade_pct);
    Serial.println("%");

    if (luminosidade_pct < 20.0f) {
        Serial.println("[ALERTA] Ambiente escuro detectado!");
    }

    delay(1000);
}
\end{codigoc}

\section{Suavizando leituras com um vetor: média móvel}

Sensores analógicos são naturalmente ruidosos — duas leituras consecutivas do mesmo ambiente raramente retornam exatamente o mesmo valor. Uma técnica simples e eficaz para suavizar esse ruído é a \textbf{média móvel}: manter um pequeno histórico das últimas N leituras (usando um vetor, como no \cref{cap:arranjos}) e trabalhar sempre com a média desse histórico, em vez de uma leitura isolada.

\begin{codigoc}{Média móvel das últimas 10 leituras}
#define PINO_SENSOR 34
#define TAMANHO_HISTORICO 10

int historico[TAMANHO_HISTORICO];
int indice_atual = 0;

void setup() {
    Serial.begin(115200);
    for (int i = 0; i < TAMANHO_HISTORICO; i++) {
        historico[i] = 0;
    }
}

void loop() {
    historico[indice_atual] = analogRead(PINO_SENSOR);
    indice_atual = (indice_atual + 1) % TAMANHO_HISTORICO;

    long soma = 0;
    for (int i = 0; i < TAMANHO_HISTORICO; i++) {
        soma += historico[i];
    }
    float media = (float) soma / TAMANHO_HISTORICO;

    Serial.print("[SENSOR] media movel = ");
    Serial.println(media);

    delay(100);
}
\end{codigoc}

O operador \code{\%} (resto da divisão, \cref{cap:tipos}) faz o índice \code{indice\_atual} voltar a 0 automaticamente ao chegar em \code{TAMANHO\_HISTORICO} — o vetor se comporta como um \textbf{buffer circular}, sempre sobrescrevendo a leitura mais antiga com a mais nova. 
